% Supplementary material for Wits NLP 2026 Project 2.
% Place this file beside paper/report.tex.
% Insert with % Supplementary material for Wits NLP 2026 Project 2.
% Place this file beside paper/report.tex.
% Insert with % Supplementary material for Wits NLP 2026 Project 2.
% Place this file beside paper/report.tex.
% Insert with % Supplementary material for Wits NLP 2026 Project 2.
% Place this file beside paper/report.tex.
% Insert with \input{supplementary} AFTER the
% references and completed, author-approved NeurIPS 2024 questionnaire,
% as specified by the course brief. Requires graphicx (in the CCN manuscript).
% This is not a replacement for the required Faculty AI ethics statement.
% This file is separate until the NeurIPS answers are author-approved.
% It must not be appended before that required questionnaire.
\clearpage
\section*{Supplementary Material}

\subsection*{S1. Sampling design and predictive evaluation}
\noindent\textbf{Corpus and separation.}
The tokenizer experiment uses the \emph{Imvo Zabantsundu} pages in the supplied
language-view archive. It trains on 1885 and January--September 1887
(342 pages, 86 publication issues), selects the count-bigram smoothing
parameter on October--December 1887 (48 pages, 12 issues), and evaluates
on 1888 (52 issues; 134 English-assigned and 208 isiXhosa-assigned
eligible page-view excerpts). The 1889--1894 material is a dependent sensitivity check (306 issues):
those years had been inspected in earlier exploratory work. Training and evaluation
issue identifiers are disjoint, but identical passages can recur in
different newspaper issues. Eligible evaluation page views contribute their first 128 alphabetic
words when at least 40 are present. English/isiXhosa views are archive
assignments, not manually adjudicated word-level languages.

\noindent\textbf{Four sampling arms.}
For eligible view $v$ of page $i$, let $n_{iv}$ be the number of ASCII
alphabetic words, $n_i=\sum_v n_{iv}$, and $N$ the number of eligible
pages. The probability of selecting a language view is either the
page-average share or the pooled word share:
\begin{equation}
\begin{aligned}
p_P(v)&=\frac{1}{N}\sum_i\frac{n_{iv}}{n_i},\\
p_W(v)&=\frac{\sum_i n_{iv}}{\sum_i n_i}.
\end{aligned}
\end{equation}
Conditional on the view, the page is drawn with probability
$q_P(i\mid v)\propto n_{iv}/n_i$ or
$q_W(i\mid v)\propto n_{iv}$, normalized over its eligible pages.
In the training pool, the expected English-view draw share is
23.55\% for $P$ and 32.56\% for $W$; realized character shares can differ
because sampled windows vary in length.
The cross of allocation ($P/W$) and within-view page selection ($P/W$)
produces $PP$, $PW$, $WP$ and $WW$, respectively; changing only the first
letter isolates the allocation factor. Each draw selects a uniformly
random eligible 32-word within-page window. The study uses lowercase
\texttt{[a-z]+} words, no stemming or modernization, and a training
budget of $500{,}000\pm30$ alphabetic characters per tokenizer.

Both BPE and Unigram use a fixed 2,000-piece vocabulary, identity
SentencePiece normalization, character coverage 1.0, single-threaded
fitting, and five sampling seeds (441--445) per arm:
$2\times4\times5=40$ main tokenizers.
The vocabulary size and tokenizer families were chosen in earlier
exploration and held fixed for this study, not selected on the 1888 set. The predictive model is a
\emph{count bigram}, not a neural language model. Each tokenizer is
applied to the same underlying 545 training page-view texts and held-out
raw excerpts; only the induced token sequence changes.

\noindent\textbf{Measures and model selection.}
Segmentation cost is $E=100P/C$, where $P$ is the number of SentencePiece
tokens and $C$ is the number of alphabetic source characters. For a token
$t$ and predecessor $h$, with token vocabulary $V$ and total training
token count $T$, the smoothed bigram probability is
\begin{equation}
\begin{aligned}
u(t)&=\frac{c(t)+0.1}{T+0.1|V|},\\
\Pr(t\mid h)&=\frac{c(h,t)+\alpha u(t)}{c(h)+\alpha}.
\end{aligned}
\end{equation}
The start-of-document predecessor is handled as a separate context.
Held-out bits per alphabetic character (BPC) is the negative base-two
log likelihood divided by $C$. The globally shared $\alpha$ was
chosen from $\{1,10,100\}$ on 344 validation windows screened for
exact training 13-word overlaps (13,760 model-window scores), averaging
view-level BPC across tokenizer families, arms and seeds. The respective validation
means were 3.1246, 2.8279 and 2.7329 BPC, selecting $\alpha=100$ before
the 1888 comparison.

For either response $Y\in\{E,\mathrm{BPC}\}$, the factorial
allocation and within-view contrasts are
\begin{equation}
\begin{aligned}
\Delta_A&=\tfrac12[(Y_{WP}-Y_{PP})+(Y_{WW}-Y_{PW})],\\
\Delta_B&=\tfrac12[(Y_{PW}-Y_{PP})+(Y_{WW}-Y_{WP})].
\end{aligned}
\end{equation}
The reported effects are means over five fitted seeds; 95\% percentile
intervals use 5,000 paired resamples of \emph{entire publication issues}
(52 issues for 1888), retaining all pages and arms of a sampled issue
together. These intervals describe uncertainty over the observed issue clusters,
conditional on the fixed trained seeds and archive; five seeds are not
five independent historical samples.

\subsection*{S2. Full tokenizer results and falsification checks}
The main paper emphasizes BPE English and the language-view trade-off.
Tables~\ref{tab:sup-seg}--\ref{tab:sup-bpc} give the complete primary
1888 allocation contrast for both models and both archive views.
Negative values indicate lower cost under word-like view allocation.

\begin{table}[ht]
\centering\footnotesize
\renewcommand{\arraystretch}{1.13}
\begin{tabular*}{\columnwidth}{@{\extracolsep{\fill}}lrr@{}}
\hline
Model/view & Effect & 95\% interval\\
\hline
BPE / EN & $-1.878$ & $[-1.972,-1.782]$\\
BPE / XH & $+0.453$ & $[+0.407,+0.497]$\\
Unigram / EN & $-1.795$ & $[-1.908,-1.671]$\\
Unigram / XH & $+0.555$ & $[+0.513,+0.600]$\\
\hline
\end{tabular*}
\caption{1888 allocation effect $\Delta_A$: subword pieces per
100 alphabetic characters, with issue-bootstrap 95\% intervals.}
\label{tab:sup-seg}
\end{table}

\begin{table}[ht]
\centering\footnotesize
\renewcommand{\arraystretch}{1.13}
\begin{tabular*}{\columnwidth}{@{\extracolsep{\fill}}lrr@{}}
\hline
Model/view & Effect & 95\% interval\\
\hline
BPE / EN & $-0.03722$ & $[-0.04236,-0.03212]$\\
BPE / XH & $-0.00275$ & $[-0.00523,-0.00014]$\\
Unigram / EN & $-0.03834$ & $[-0.04366,-0.03318]$\\
Unigram / XH & $+0.01048$ & $[+0.00805,+0.01286]$\\
\hline
\end{tabular*}
\caption{1888 $\Delta_A$: count-bigram bits per alphabetic character.
Held-out source texts are matched across arms.}
\label{tab:sup-bpc}
\end{table}

BPE's XH bigram change is small and has mixed signs across seeds
(three negative, two positive), unlike its five-of-five positive
segmentation effect. Unigram has higher XH segmentation \emph{and}
higher XH bigram loss. Thus a tokenizer can encode the archive more
compactly in one view without consistently improving predictive loss
in the other. In English-view BPE, the separately estimated
within-view selection effect on BPC is $+0.01938$
(95\% interval $[+0.01146,+0.02769]$), opposing the allocation benefit.

For the repetition-control experiment, 20 additional BPE tokenizers
were trained under $PP/WW$, five seeds and two matched-support
conditions. The 1888 English $E_{PP}-E_{WW}$ contrast after filtering
training-only repeated 13-word sequences was $+2.389$
$[+2.182,+2.590]$ pieces per 100 characters; an unfiltered model on
the same page-view support gave $+2.097$ $[+1.905,+2.280]$.
Their paired difference is $+0.292$ $[+0.163,+0.421]$.
The trade-off survives this exact reprint filter;
the filter cannot rule out shorter, paraphrased or otherwise shared
historical passages.

A counterexample concerns the \emph{whole-policy}, rather than
allocation-only, contrast. On dependent 1889--1894 English data,
BPE $\mathrm{BPC}_{WW}-\mathrm{BPC}_{PP}=+0.00450$; the issue-bootstrap
interval is $[+0.00176,+0.00728]$, but the six-year resampling
interval $[-0.00268,+0.01305]$ includes zero. This does not support
a reliable improvement from changing the entire sampling policy.

\subsection*{S3. ArchiveLens estimands and statistical calibration}
ArchiveLens uses the archive's XH-assigned views from both
\emph{Imvo Zabantsundu} and \emph{Isigidimi sama-Xosa}.
It fixes the 120 most frequent within-word four-character sequences
by the mean of the two newspapers' normalized \emph{1887} rates;
there is no selection from 1885/1888 outcome differences.
A page is eligible if it has at least 40 alphabetic words.
The 1885/1888 support comprises 195/208 pages (49/52 issues)
for \emph{Imvo} and 95/96 pages (12/12 issues) for \emph{Isigidimi}.

For pattern $g$ on page $i$, let $c_{ig}$ be its occurrence count and
$d_i=\sum_{w\in i}\max(0,|w|-3)$ the number of eligible within-word
four-gram positions. The \emph{mass} rate pools all counts and
positions, while \emph{page} gives each eligible page the same weight.
The \emph{issue} rate first pools pages within an issue, then averages
issue rates. The \emph{month} rate averages issue rates within each
calendar month before weighting the 12 shared calendar months equally.
All rates are multiplied by $10^4$. Each reported historical contrast
is the 1888 rate minus the 1885 rate. These target different exposure units rather than a single,
unambiguously correct historical population.

A 3,000-draw nested month/issue bootstrap produces descriptive
95\% intervals, not multiplicity-adjusted discovery claims.
Separately, 4,999 year-label permutations within calendar months move
whole publication issues, preserve month/year issue counts and retain
page groups. For each pattern, one test examines
$|\Delta_{g,\mathrm{mass}}|$ (temporal change), while another examines
$\max_{w\ne\mathrm{mass}}|\Delta_{g,w}-\Delta_{g,\mathrm{mass}}|$
(estimator sensitivity). Each family of 120 pattern tests within a
newspaper receives its own Holm correction. A third benchmark examines
the \emph{number} of sign reversals, not the reliability of each
individual reversal warning.

\begin{table}[H]
\centering\footnotesize
\renewcommand{\arraystretch}{1.13}
\begin{tabular*}{\columnwidth}{@{\extracolsep{\fill}}lrr@{}}
\hline
ArchiveLens result & Imvo & Isigidimi\\
\hline
Reversal warnings & 23/120 & 4/120\\
Null mean & 29.02 & 8.30\\
Central 95\% null interval & 19--40 & 3--16\\
Upper-tail $p$ & 0.8882 & 0.9486\\
Holm temporal effects & 26 & 0\\
Holm estimator gaps & 29 & 0\\
\hline
\end{tabular*}
\caption{Reversal-count calibration and separately Holm-adjusted year and estimator-gap tests.}
\label{tab:sup-arch}
\end{table}

For context, the earlier \emph{unadjusted} bootstrap intervals
excluded zero for 38 \emph{Imvo} and 21 \emph{Isigidimi}
mass-weighted year changes. Those descriptive interval counts are not
interchangeable with the Holm-corrected permutation tests.

Neither newspaper has an unexpectedly high reversal \emph{count}
under this conditional reference. The 29 adjusted \emph{Imvo}
estimator gaps answer a different question: nine also had original
sign-reversal warnings. Table~\ref{tab:sup-examples} shows several
of those descriptive disagreements, using the same patterns and
original corpus. For \texttt{ants}, the maximum gap has
Holm-adjusted $p=0.024$, whereas the \emph{mass-weighted temporal}
test has adjusted $p=1$. The gap is evidence of estimator sensitivity,
not a pair of proven historical trends.

\begin{table}[H]
\centering\footnotesize
\renewcommand{\arraystretch}{1.15}
\begin{tabular*}{\columnwidth}{@{\extracolsep{\fill}}lrrrr@{}}
\hline
Fragment & Mass & Page & Issue & Month\\
\hline
\texttt{ants} & $-0.23$ & $+4.80$ & $-0.03$ & $+0.15$\\
\texttt{ntsu} & $-2.03$ & $+0.81$ & $-1.79$ & $-1.65$\\
\texttt{tsha} & $+0.40$ & $-2.23$ & $+0.20$ & $+0.12$\\
\texttt{teto} & $-0.36$ & $+2.21$ & $-0.12$ & $-0.29$\\
\hline
\end{tabular*}
\caption{Illustrative \emph{Imvo} 1888-minus-1885 rate
changes per 10,000 eligible four-gram positions. Patterns are
surface fragments, not validated isiXhosa words or meanings;
individual sign changes are not adjusted discoveries.}
\label{tab:sup-examples}
\end{table}

\clearpage
\subsection*{S4. Archive structure, repeated text, and validity}
Figure~\ref{fig:sup-page} shows how eligible page text varies with
position. \emph{Imvo} page 2 has 4.24 times the within-word four-gram
opportunities of page 1 in 1885, and 3.69 times in 1888; the
\emph{Isigidimi} profile is considerably flatter. The same pattern
can therefore carry different weight in character-mass and equal-page
comparisons. Page position is not an annotated genre or evidence of
editorial motives.

\begin{figure}[H]
\centering
\includegraphics[width=.99\columnwidth]{supp_page_position_compact.pdf}
\caption{Mean eligible four-gram positions per page relative to
position 1 in each newspaper and year, using the original XH-assigned
archive text. Positions 1--4 are present in all four source/year
combinations. This is exposure, not a rate of a selected pattern.}
\label{fig:sup-page}
\end{figure}

The follow-up masks words covered by exact 13-word sequences
repeated across \emph{different} issues within the same newspaper.
Table~\ref{tab:sup-cover} shows the text affected. After masking,
ten original \emph{Imvo} reversal warnings and one
\emph{Isigidimi} warning disappear;
this changes the eligible text and is a sensitivity analysis, not an
estimate of the causal effect of ``removing duplication.'' Among the
32 originally flagged \emph{Imvo} patterns, seven can change their
mass-weighted year-difference sign when one complete issue is omitted;
four of six flagged \emph{Isigidimi} patterns do so.

\begin{table}[H]
\centering\footnotesize
\renewcommand{\arraystretch}{1.15}
\begin{tabular*}{\columnwidth}{@{\extracolsep{\fill}}lrrr@{}}
\hline
Newspaper / year & Pages & Issues & Masked words\\
\hline
\emph{Imvo} 1885 & 195 & 49 & 20.0\%\\
\emph{Imvo} 1888 & 208 & 52 & 21.9\%\\
\emph{Isigidimi} 1885 & 95 & 12 & 2.4\%\\
\emph{Isigidimi} 1888 & 96 & 12 & 1.3\%\\
\hline
\end{tabular*}
\caption{Original eligible page support and fraction of words
covered by repeated passages in each newspaper and year.
The mask also removes some pages from the sensitivity analysis.}
\label{tab:sup-cover}
\end{table}

The tokenizer experiment also has disjoint date-based splits but not
perfectly novel evaluation \emph{content}: 514 of 2,493 held-out
page-view excerpts (20.6\%) contain an exact 13-word sequence also
seen in training-era text. Restricting to 99 nonoverlapping 1888 English-view excerpts yields a
post-hoc BPE bigram allocation effect of $-0.02874$ BPC, compared with
$-0.03722$ on full support. The direction remains, but the restriction
changes the evaluation population, and shorter or paraphrased overlaps
remain possible.

ArchiveLens uses the original text views, not fitted tokenizer
outputs. The two studies share a question about measurement, not a
measured causal pathway. Both use surviving transcriptions from one
archive, with no experimental control of editorial composition,
preservation or OCR/transcription effects. The historical years
had been examined before the final follow-ups, and permutation
exchangeability is a conditional reference rather than a credible
assumption of identical historical eras. The findings support
measurement- and sampling-sensitivity diagnostics, not claims about
historical authors' beliefs, semantic change, social discrimination,
or downstream neural-model fairness.

\noindent\textbf{Generative-AI assistance.}
ChatGPT assisted experiment planning, analysis/code checks, figure preparation,
and report drafting/editing; NotebookLM and Gemini assisted literature
exploration and citations, as declared. The authors remain responsible
for validating the code, numerical evidence and references.

\noindent\textbf{Code and reproduction.}
The recorded model and ArchiveLens results are verified by the
repository's Python 3.13 checks; the full tokenizer retraining requires
the original source ZIP and is not run in routine CI. The source archive
SHA-256 is \texttt{375c5864...be321d5} (full value in
\texttt{README.txt}). Hardware and runtime measurements were not
retained consistently, so none are estimated here. Archive redistribution
rights require separate confirmation.
The tables can be checked in
\path{results/segmentation.csv}, \path{results/bigram.csv}, and
\path{results/repetition.csv}. The 120-pattern estimands and their
adjusted calibration tests are in \path{results/archivelens/}; numerical inputs for
Fig.~\ref{fig:sup-page} are generated directly from the original ZIP
by \path{paper/analysis/make_supp_figure.py}; the independently
recounted inputs are in \path{paper/analysis/supp_page_position_summary.csv}.
Reproduction commands appear in the required \path{README.txt}.
 AFTER the
% references and completed, author-approved NeurIPS 2024 questionnaire,
% as specified by the course brief. Requires graphicx (in the CCN manuscript).
% This is not a replacement for the required Faculty AI ethics statement.
% This file is separate until the NeurIPS answers are author-approved.
% It must not be appended before that required questionnaire.
\clearpage
\section*{Supplementary Material}

\subsection*{S1. Sampling design and predictive evaluation}
\noindent\textbf{Corpus and separation.}
The tokenizer experiment uses the \emph{Imvo Zabantsundu} pages in the supplied
language-view archive. It trains on 1885 and January--September 1887
(342 pages, 86 publication issues), selects the count-bigram smoothing
parameter on October--December 1887 (48 pages, 12 issues), and evaluates
on 1888 (52 issues; 134 English-assigned and 208 isiXhosa-assigned
eligible page-view excerpts). The 1889--1894 material is a dependent sensitivity check (306 issues):
those years had been inspected in earlier exploratory work. Training and evaluation
issue identifiers are disjoint, but identical passages can recur in
different newspaper issues. Eligible evaluation page views contribute their first 128 alphabetic
words when at least 40 are present. English/isiXhosa views are archive
assignments, not manually adjudicated word-level languages.

\noindent\textbf{Four sampling arms.}
For eligible view $v$ of page $i$, let $n_{iv}$ be the number of ASCII
alphabetic words, $n_i=\sum_v n_{iv}$, and $N$ the number of eligible
pages. The probability of selecting a language view is either the
page-average share or the pooled word share:
\begin{equation}
\begin{aligned}
p_P(v)&=\frac{1}{N}\sum_i\frac{n_{iv}}{n_i},\\
p_W(v)&=\frac{\sum_i n_{iv}}{\sum_i n_i}.
\end{aligned}
\end{equation}
Conditional on the view, the page is drawn with probability
$q_P(i\mid v)\propto n_{iv}/n_i$ or
$q_W(i\mid v)\propto n_{iv}$, normalized over its eligible pages.
In the training pool, the expected English-view draw share is
23.55\% for $P$ and 32.56\% for $W$; realized character shares can differ
because sampled windows vary in length.
The cross of allocation ($P/W$) and within-view page selection ($P/W$)
produces $PP$, $PW$, $WP$ and $WW$, respectively; changing only the first
letter isolates the allocation factor. Each draw selects a uniformly
random eligible 32-word within-page window. The study uses lowercase
\texttt{[a-z]+} words, no stemming or modernization, and a training
budget of $500{,}000\pm30$ alphabetic characters per tokenizer.

Both BPE and Unigram use a fixed 2,000-piece vocabulary, identity
SentencePiece normalization, character coverage 1.0, single-threaded
fitting, and five sampling seeds (441--445) per arm:
$2\times4\times5=40$ main tokenizers.
The vocabulary size and tokenizer families were chosen in earlier
exploration and held fixed for this study, not selected on the 1888 set. The predictive model is a
\emph{count bigram}, not a neural language model. Each tokenizer is
applied to the same underlying 545 training page-view texts and held-out
raw excerpts; only the induced token sequence changes.

\noindent\textbf{Measures and model selection.}
Segmentation cost is $E=100P/C$, where $P$ is the number of SentencePiece
tokens and $C$ is the number of alphabetic source characters. For a token
$t$ and predecessor $h$, with token vocabulary $V$ and total training
token count $T$, the smoothed bigram probability is
\begin{equation}
\begin{aligned}
u(t)&=\frac{c(t)+0.1}{T+0.1|V|},\\
\Pr(t\mid h)&=\frac{c(h,t)+\alpha u(t)}{c(h)+\alpha}.
\end{aligned}
\end{equation}
The start-of-document predecessor is handled as a separate context.
Held-out bits per alphabetic character (BPC) is the negative base-two
log likelihood divided by $C$. The globally shared $\alpha$ was
chosen from $\{1,10,100\}$ on 344 validation windows screened for
exact training 13-word overlaps (13,760 model-window scores), averaging
view-level BPC across tokenizer families, arms and seeds. The respective validation
means were 3.1246, 2.8279 and 2.7329 BPC, selecting $\alpha=100$ before
the 1888 comparison.

For either response $Y\in\{E,\mathrm{BPC}\}$, the factorial
allocation and within-view contrasts are
\begin{equation}
\begin{aligned}
\Delta_A&=\tfrac12[(Y_{WP}-Y_{PP})+(Y_{WW}-Y_{PW})],\\
\Delta_B&=\tfrac12[(Y_{PW}-Y_{PP})+(Y_{WW}-Y_{WP})].
\end{aligned}
\end{equation}
The reported effects are means over five fitted seeds; 95\% percentile
intervals use 5,000 paired resamples of \emph{entire publication issues}
(52 issues for 1888), retaining all pages and arms of a sampled issue
together. These intervals describe uncertainty over the observed issue clusters,
conditional on the fixed trained seeds and archive; five seeds are not
five independent historical samples.

\subsection*{S2. Full tokenizer results and falsification checks}
The main paper emphasizes BPE English and the language-view trade-off.
Tables~\ref{tab:sup-seg}--\ref{tab:sup-bpc} give the complete primary
1888 allocation contrast for both models and both archive views.
Negative values indicate lower cost under word-like view allocation.

\begin{table}[ht]
\centering\footnotesize
\renewcommand{\arraystretch}{1.13}
\begin{tabular*}{\columnwidth}{@{\extracolsep{\fill}}lrr@{}}
\hline
Model/view & Effect & 95\% interval\\
\hline
BPE / EN & $-1.878$ & $[-1.972,-1.782]$\\
BPE / XH & $+0.453$ & $[+0.407,+0.497]$\\
Unigram / EN & $-1.795$ & $[-1.908,-1.671]$\\
Unigram / XH & $+0.555$ & $[+0.513,+0.600]$\\
\hline
\end{tabular*}
\caption{1888 allocation effect $\Delta_A$: subword pieces per
100 alphabetic characters, with issue-bootstrap 95\% intervals.}
\label{tab:sup-seg}
\end{table}

\begin{table}[ht]
\centering\footnotesize
\renewcommand{\arraystretch}{1.13}
\begin{tabular*}{\columnwidth}{@{\extracolsep{\fill}}lrr@{}}
\hline
Model/view & Effect & 95\% interval\\
\hline
BPE / EN & $-0.03722$ & $[-0.04236,-0.03212]$\\
BPE / XH & $-0.00275$ & $[-0.00523,-0.00014]$\\
Unigram / EN & $-0.03834$ & $[-0.04366,-0.03318]$\\
Unigram / XH & $+0.01048$ & $[+0.00805,+0.01286]$\\
\hline
\end{tabular*}
\caption{1888 $\Delta_A$: count-bigram bits per alphabetic character.
Held-out source texts are matched across arms.}
\label{tab:sup-bpc}
\end{table}

BPE's XH bigram change is small and has mixed signs across seeds
(three negative, two positive), unlike its five-of-five positive
segmentation effect. Unigram has higher XH segmentation \emph{and}
higher XH bigram loss. Thus a tokenizer can encode the archive more
compactly in one view without consistently improving predictive loss
in the other. In English-view BPE, the separately estimated
within-view selection effect on BPC is $+0.01938$
(95\% interval $[+0.01146,+0.02769]$), opposing the allocation benefit.

For the repetition-control experiment, 20 additional BPE tokenizers
were trained under $PP/WW$, five seeds and two matched-support
conditions. The 1888 English $E_{PP}-E_{WW}$ contrast after filtering
training-only repeated 13-word sequences was $+2.389$
$[+2.182,+2.590]$ pieces per 100 characters; an unfiltered model on
the same page-view support gave $+2.097$ $[+1.905,+2.280]$.
Their paired difference is $+0.292$ $[+0.163,+0.421]$.
The trade-off survives this exact reprint filter;
the filter cannot rule out shorter, paraphrased or otherwise shared
historical passages.

A counterexample concerns the \emph{whole-policy}, rather than
allocation-only, contrast. On dependent 1889--1894 English data,
BPE $\mathrm{BPC}_{WW}-\mathrm{BPC}_{PP}=+0.00450$; the issue-bootstrap
interval is $[+0.00176,+0.00728]$, but the six-year resampling
interval $[-0.00268,+0.01305]$ includes zero. This does not support
a reliable improvement from changing the entire sampling policy.

\subsection*{S3. ArchiveLens estimands and statistical calibration}
ArchiveLens uses the archive's XH-assigned views from both
\emph{Imvo Zabantsundu} and \emph{Isigidimi sama-Xosa}.
It fixes the 120 most frequent within-word four-character sequences
by the mean of the two newspapers' normalized \emph{1887} rates;
there is no selection from 1885/1888 outcome differences.
A page is eligible if it has at least 40 alphabetic words.
The 1885/1888 support comprises 195/208 pages (49/52 issues)
for \emph{Imvo} and 95/96 pages (12/12 issues) for \emph{Isigidimi}.

For pattern $g$ on page $i$, let $c_{ig}$ be its occurrence count and
$d_i=\sum_{w\in i}\max(0,|w|-3)$ the number of eligible within-word
four-gram positions. The \emph{mass} rate pools all counts and
positions, while \emph{page} gives each eligible page the same weight.
The \emph{issue} rate first pools pages within an issue, then averages
issue rates. The \emph{month} rate averages issue rates within each
calendar month before weighting the 12 shared calendar months equally.
All rates are multiplied by $10^4$. Each reported historical contrast
is the 1888 rate minus the 1885 rate. These target different exposure units rather than a single,
unambiguously correct historical population.

A 3,000-draw nested month/issue bootstrap produces descriptive
95\% intervals, not multiplicity-adjusted discovery claims.
Separately, 4,999 year-label permutations within calendar months move
whole publication issues, preserve month/year issue counts and retain
page groups. For each pattern, one test examines
$|\Delta_{g,\mathrm{mass}}|$ (temporal change), while another examines
$\max_{w\ne\mathrm{mass}}|\Delta_{g,w}-\Delta_{g,\mathrm{mass}}|$
(estimator sensitivity). Each family of 120 pattern tests within a
newspaper receives its own Holm correction. A third benchmark examines
the \emph{number} of sign reversals, not the reliability of each
individual reversal warning.

\begin{table}[H]
\centering\footnotesize
\renewcommand{\arraystretch}{1.13}
\begin{tabular*}{\columnwidth}{@{\extracolsep{\fill}}lrr@{}}
\hline
ArchiveLens result & Imvo & Isigidimi\\
\hline
Reversal warnings & 23/120 & 4/120\\
Null mean & 29.02 & 8.30\\
Central 95\% null interval & 19--40 & 3--16\\
Upper-tail $p$ & 0.8882 & 0.9486\\
Holm temporal effects & 26 & 0\\
Holm estimator gaps & 29 & 0\\
\hline
\end{tabular*}
\caption{Reversal-count calibration and separately Holm-adjusted year and estimator-gap tests.}
\label{tab:sup-arch}
\end{table}

For context, the earlier \emph{unadjusted} bootstrap intervals
excluded zero for 38 \emph{Imvo} and 21 \emph{Isigidimi}
mass-weighted year changes. Those descriptive interval counts are not
interchangeable with the Holm-corrected permutation tests.

Neither newspaper has an unexpectedly high reversal \emph{count}
under this conditional reference. The 29 adjusted \emph{Imvo}
estimator gaps answer a different question: nine also had original
sign-reversal warnings. Table~\ref{tab:sup-examples} shows several
of those descriptive disagreements, using the same patterns and
original corpus. For \texttt{ants}, the maximum gap has
Holm-adjusted $p=0.024$, whereas the \emph{mass-weighted temporal}
test has adjusted $p=1$. The gap is evidence of estimator sensitivity,
not a pair of proven historical trends.

\begin{table}[H]
\centering\footnotesize
\renewcommand{\arraystretch}{1.15}
\begin{tabular*}{\columnwidth}{@{\extracolsep{\fill}}lrrrr@{}}
\hline
Fragment & Mass & Page & Issue & Month\\
\hline
\texttt{ants} & $-0.23$ & $+4.80$ & $-0.03$ & $+0.15$\\
\texttt{ntsu} & $-2.03$ & $+0.81$ & $-1.79$ & $-1.65$\\
\texttt{tsha} & $+0.40$ & $-2.23$ & $+0.20$ & $+0.12$\\
\texttt{teto} & $-0.36$ & $+2.21$ & $-0.12$ & $-0.29$\\
\hline
\end{tabular*}
\caption{Illustrative \emph{Imvo} 1888-minus-1885 rate
changes per 10,000 eligible four-gram positions. Patterns are
surface fragments, not validated isiXhosa words or meanings;
individual sign changes are not adjusted discoveries.}
\label{tab:sup-examples}
\end{table}

\clearpage
\subsection*{S4. Archive structure, repeated text, and validity}
Figure~\ref{fig:sup-page} shows how eligible page text varies with
position. \emph{Imvo} page 2 has 4.24 times the within-word four-gram
opportunities of page 1 in 1885, and 3.69 times in 1888; the
\emph{Isigidimi} profile is considerably flatter. The same pattern
can therefore carry different weight in character-mass and equal-page
comparisons. Page position is not an annotated genre or evidence of
editorial motives.

\begin{figure}[H]
\centering
\includegraphics[width=.99\columnwidth]{supp_page_position_compact.pdf}
\caption{Mean eligible four-gram positions per page relative to
position 1 in each newspaper and year, using the original XH-assigned
archive text. Positions 1--4 are present in all four source/year
combinations. This is exposure, not a rate of a selected pattern.}
\label{fig:sup-page}
\end{figure}

The follow-up masks words covered by exact 13-word sequences
repeated across \emph{different} issues within the same newspaper.
Table~\ref{tab:sup-cover} shows the text affected. After masking,
ten original \emph{Imvo} reversal warnings and one
\emph{Isigidimi} warning disappear;
this changes the eligible text and is a sensitivity analysis, not an
estimate of the causal effect of ``removing duplication.'' Among the
32 originally flagged \emph{Imvo} patterns, seven can change their
mass-weighted year-difference sign when one complete issue is omitted;
four of six flagged \emph{Isigidimi} patterns do so.

\begin{table}[H]
\centering\footnotesize
\renewcommand{\arraystretch}{1.15}
\begin{tabular*}{\columnwidth}{@{\extracolsep{\fill}}lrrr@{}}
\hline
Newspaper / year & Pages & Issues & Masked words\\
\hline
\emph{Imvo} 1885 & 195 & 49 & 20.0\%\\
\emph{Imvo} 1888 & 208 & 52 & 21.9\%\\
\emph{Isigidimi} 1885 & 95 & 12 & 2.4\%\\
\emph{Isigidimi} 1888 & 96 & 12 & 1.3\%\\
\hline
\end{tabular*}
\caption{Original eligible page support and fraction of words
covered by repeated passages in each newspaper and year.
The mask also removes some pages from the sensitivity analysis.}
\label{tab:sup-cover}
\end{table}

The tokenizer experiment also has disjoint date-based splits but not
perfectly novel evaluation \emph{content}: 514 of 2,493 held-out
page-view excerpts (20.6\%) contain an exact 13-word sequence also
seen in training-era text. Restricting to 99 nonoverlapping 1888 English-view excerpts yields a
post-hoc BPE bigram allocation effect of $-0.02874$ BPC, compared with
$-0.03722$ on full support. The direction remains, but the restriction
changes the evaluation population, and shorter or paraphrased overlaps
remain possible.

ArchiveLens uses the original text views, not fitted tokenizer
outputs. The two studies share a question about measurement, not a
measured causal pathway. Both use surviving transcriptions from one
archive, with no experimental control of editorial composition,
preservation or OCR/transcription effects. The historical years
had been examined before the final follow-ups, and permutation
exchangeability is a conditional reference rather than a credible
assumption of identical historical eras. The findings support
measurement- and sampling-sensitivity diagnostics, not claims about
historical authors' beliefs, semantic change, social discrimination,
or downstream neural-model fairness.

\noindent\textbf{Generative-AI assistance.}
ChatGPT assisted experiment planning, analysis/code checks, figure preparation,
and report drafting/editing; NotebookLM and Gemini assisted literature
exploration and citations, as declared. The authors remain responsible
for validating the code, numerical evidence and references.

\noindent\textbf{Code and reproduction.}
The recorded model and ArchiveLens results are verified by the
repository's Python 3.13 checks; the full tokenizer retraining requires
the original source ZIP and is not run in routine CI. The source archive
SHA-256 is \texttt{375c5864...be321d5} (full value in
\texttt{README.txt}). Hardware and runtime measurements were not
retained consistently, so none are estimated here. Archive redistribution
rights require separate confirmation.
The tables can be checked in
\path{results/segmentation.csv}, \path{results/bigram.csv}, and
\path{results/repetition.csv}. The 120-pattern estimands and their
adjusted calibration tests are in \path{results/archivelens/}; numerical inputs for
Fig.~\ref{fig:sup-page} are generated directly from the original ZIP
by \path{paper/analysis/make_supp_figure.py}; the independently
recounted inputs are in \path{paper/analysis/supp_page_position_summary.csv}.
Reproduction commands appear in the required \path{README.txt}.
 AFTER the
% references and completed, author-approved NeurIPS 2024 questionnaire,
% as specified by the course brief. Requires graphicx (in the CCN manuscript).
% This is not a replacement for the required Faculty AI ethics statement.
% This file is separate until the NeurIPS answers are author-approved.
% It must not be appended before that required questionnaire.
\clearpage
\section*{Supplementary Material}

\subsection*{S1. Sampling design and predictive evaluation}
\noindent\textbf{Corpus and separation.}
The tokenizer experiment uses the \emph{Imvo Zabantsundu} pages in the supplied
language-view archive. It trains on 1885 and January--September 1887
(342 pages, 86 publication issues), selects the count-bigram smoothing
parameter on October--December 1887 (48 pages, 12 issues), and evaluates
on 1888 (52 issues; 134 English-assigned and 208 isiXhosa-assigned
eligible page-view excerpts). The 1889--1894 material is a dependent sensitivity check (306 issues):
those years had been inspected in earlier exploratory work. Training and evaluation
issue identifiers are disjoint, but identical passages can recur in
different newspaper issues. Eligible evaluation page views contribute their first 128 alphabetic
words when at least 40 are present. English/isiXhosa views are archive
assignments, not manually adjudicated word-level languages.

\noindent\textbf{Four sampling arms.}
For eligible view $v$ of page $i$, let $n_{iv}$ be the number of ASCII
alphabetic words, $n_i=\sum_v n_{iv}$, and $N$ the number of eligible
pages. The probability of selecting a language view is either the
page-average share or the pooled word share:
\begin{equation}
\begin{aligned}
p_P(v)&=\frac{1}{N}\sum_i\frac{n_{iv}}{n_i},\\
p_W(v)&=\frac{\sum_i n_{iv}}{\sum_i n_i}.
\end{aligned}
\end{equation}
Conditional on the view, the page is drawn with probability
$q_P(i\mid v)\propto n_{iv}/n_i$ or
$q_W(i\mid v)\propto n_{iv}$, normalized over its eligible pages.
In the training pool, the expected English-view draw share is
23.55\% for $P$ and 32.56\% for $W$; realized character shares can differ
because sampled windows vary in length.
The cross of allocation ($P/W$) and within-view page selection ($P/W$)
produces $PP$, $PW$, $WP$ and $WW$, respectively; changing only the first
letter isolates the allocation factor. Each draw selects a uniformly
random eligible 32-word within-page window. The study uses lowercase
\texttt{[a-z]+} words, no stemming or modernization, and a training
budget of $500{,}000\pm30$ alphabetic characters per tokenizer.

Both BPE and Unigram use a fixed 2,000-piece vocabulary, identity
SentencePiece normalization, character coverage 1.0, single-threaded
fitting, and five sampling seeds (441--445) per arm:
$2\times4\times5=40$ main tokenizers.
The vocabulary size and tokenizer families were chosen in earlier
exploration and held fixed for this study, not selected on the 1888 set. The predictive model is a
\emph{count bigram}, not a neural language model. Each tokenizer is
applied to the same underlying 545 training page-view texts and held-out
raw excerpts; only the induced token sequence changes.

\noindent\textbf{Measures and model selection.}
Segmentation cost is $E=100P/C$, where $P$ is the number of SentencePiece
tokens and $C$ is the number of alphabetic source characters. For a token
$t$ and predecessor $h$, with token vocabulary $V$ and total training
token count $T$, the smoothed bigram probability is
\begin{equation}
\begin{aligned}
u(t)&=\frac{c(t)+0.1}{T+0.1|V|},\\
\Pr(t\mid h)&=\frac{c(h,t)+\alpha u(t)}{c(h)+\alpha}.
\end{aligned}
\end{equation}
The start-of-document predecessor is handled as a separate context.
Held-out bits per alphabetic character (BPC) is the negative base-two
log likelihood divided by $C$. The globally shared $\alpha$ was
chosen from $\{1,10,100\}$ on 344 validation windows screened for
exact training 13-word overlaps (13,760 model-window scores), averaging
view-level BPC across tokenizer families, arms and seeds. The respective validation
means were 3.1246, 2.8279 and 2.7329 BPC, selecting $\alpha=100$ before
the 1888 comparison.

For either response $Y\in\{E,\mathrm{BPC}\}$, the factorial
allocation and within-view contrasts are
\begin{equation}
\begin{aligned}
\Delta_A&=\tfrac12[(Y_{WP}-Y_{PP})+(Y_{WW}-Y_{PW})],\\
\Delta_B&=\tfrac12[(Y_{PW}-Y_{PP})+(Y_{WW}-Y_{WP})].
\end{aligned}
\end{equation}
The reported effects are means over five fitted seeds; 95\% percentile
intervals use 5,000 paired resamples of \emph{entire publication issues}
(52 issues for 1888), retaining all pages and arms of a sampled issue
together. These intervals describe uncertainty over the observed issue clusters,
conditional on the fixed trained seeds and archive; five seeds are not
five independent historical samples.

\subsection*{S2. Full tokenizer results and falsification checks}
The main paper emphasizes BPE English and the language-view trade-off.
Tables~\ref{tab:sup-seg}--\ref{tab:sup-bpc} give the complete primary
1888 allocation contrast for both models and both archive views.
Negative values indicate lower cost under word-like view allocation.

\begin{table}[ht]
\centering\footnotesize
\renewcommand{\arraystretch}{1.13}
\begin{tabular*}{\columnwidth}{@{\extracolsep{\fill}}lrr@{}}
\hline
Model/view & Effect & 95\% interval\\
\hline
BPE / EN & $-1.878$ & $[-1.972,-1.782]$\\
BPE / XH & $+0.453$ & $[+0.407,+0.497]$\\
Unigram / EN & $-1.795$ & $[-1.908,-1.671]$\\
Unigram / XH & $+0.555$ & $[+0.513,+0.600]$\\
\hline
\end{tabular*}
\caption{1888 allocation effect $\Delta_A$: subword pieces per
100 alphabetic characters, with issue-bootstrap 95\% intervals.}
\label{tab:sup-seg}
\end{table}

\begin{table}[ht]
\centering\footnotesize
\renewcommand{\arraystretch}{1.13}
\begin{tabular*}{\columnwidth}{@{\extracolsep{\fill}}lrr@{}}
\hline
Model/view & Effect & 95\% interval\\
\hline
BPE / EN & $-0.03722$ & $[-0.04236,-0.03212]$\\
BPE / XH & $-0.00275$ & $[-0.00523,-0.00014]$\\
Unigram / EN & $-0.03834$ & $[-0.04366,-0.03318]$\\
Unigram / XH & $+0.01048$ & $[+0.00805,+0.01286]$\\
\hline
\end{tabular*}
\caption{1888 $\Delta_A$: count-bigram bits per alphabetic character.
Held-out source texts are matched across arms.}
\label{tab:sup-bpc}
\end{table}

BPE's XH bigram change is small and has mixed signs across seeds
(three negative, two positive), unlike its five-of-five positive
segmentation effect. Unigram has higher XH segmentation \emph{and}
higher XH bigram loss. Thus a tokenizer can encode the archive more
compactly in one view without consistently improving predictive loss
in the other. In English-view BPE, the separately estimated
within-view selection effect on BPC is $+0.01938$
(95\% interval $[+0.01146,+0.02769]$), opposing the allocation benefit.

For the repetition-control experiment, 20 additional BPE tokenizers
were trained under $PP/WW$, five seeds and two matched-support
conditions. The 1888 English $E_{PP}-E_{WW}$ contrast after filtering
training-only repeated 13-word sequences was $+2.389$
$[+2.182,+2.590]$ pieces per 100 characters; an unfiltered model on
the same page-view support gave $+2.097$ $[+1.905,+2.280]$.
Their paired difference is $+0.292$ $[+0.163,+0.421]$.
The trade-off survives this exact reprint filter;
the filter cannot rule out shorter, paraphrased or otherwise shared
historical passages.

A counterexample concerns the \emph{whole-policy}, rather than
allocation-only, contrast. On dependent 1889--1894 English data,
BPE $\mathrm{BPC}_{WW}-\mathrm{BPC}_{PP}=+0.00450$; the issue-bootstrap
interval is $[+0.00176,+0.00728]$, but the six-year resampling
interval $[-0.00268,+0.01305]$ includes zero. This does not support
a reliable improvement from changing the entire sampling policy.

\subsection*{S3. ArchiveLens estimands and statistical calibration}
ArchiveLens uses the archive's XH-assigned views from both
\emph{Imvo Zabantsundu} and \emph{Isigidimi sama-Xosa}.
It fixes the 120 most frequent within-word four-character sequences
by the mean of the two newspapers' normalized \emph{1887} rates;
there is no selection from 1885/1888 outcome differences.
A page is eligible if it has at least 40 alphabetic words.
The 1885/1888 support comprises 195/208 pages (49/52 issues)
for \emph{Imvo} and 95/96 pages (12/12 issues) for \emph{Isigidimi}.

For pattern $g$ on page $i$, let $c_{ig}$ be its occurrence count and
$d_i=\sum_{w\in i}\max(0,|w|-3)$ the number of eligible within-word
four-gram positions. The \emph{mass} rate pools all counts and
positions, while \emph{page} gives each eligible page the same weight.
The \emph{issue} rate first pools pages within an issue, then averages
issue rates. The \emph{month} rate averages issue rates within each
calendar month before weighting the 12 shared calendar months equally.
All rates are multiplied by $10^4$. Each reported historical contrast
is the 1888 rate minus the 1885 rate. These target different exposure units rather than a single,
unambiguously correct historical population.

A 3,000-draw nested month/issue bootstrap produces descriptive
95\% intervals, not multiplicity-adjusted discovery claims.
Separately, 4,999 year-label permutations within calendar months move
whole publication issues, preserve month/year issue counts and retain
page groups. For each pattern, one test examines
$|\Delta_{g,\mathrm{mass}}|$ (temporal change), while another examines
$\max_{w\ne\mathrm{mass}}|\Delta_{g,w}-\Delta_{g,\mathrm{mass}}|$
(estimator sensitivity). Each family of 120 pattern tests within a
newspaper receives its own Holm correction. A third benchmark examines
the \emph{number} of sign reversals, not the reliability of each
individual reversal warning.

\begin{table}[H]
\centering\footnotesize
\renewcommand{\arraystretch}{1.13}
\begin{tabular*}{\columnwidth}{@{\extracolsep{\fill}}lrr@{}}
\hline
ArchiveLens result & Imvo & Isigidimi\\
\hline
Reversal warnings & 23/120 & 4/120\\
Null mean & 29.02 & 8.30\\
Central 95\% null interval & 19--40 & 3--16\\
Upper-tail $p$ & 0.8882 & 0.9486\\
Holm temporal effects & 26 & 0\\
Holm estimator gaps & 29 & 0\\
\hline
\end{tabular*}
\caption{Reversal-count calibration and separately Holm-adjusted year and estimator-gap tests.}
\label{tab:sup-arch}
\end{table}

For context, the earlier \emph{unadjusted} bootstrap intervals
excluded zero for 38 \emph{Imvo} and 21 \emph{Isigidimi}
mass-weighted year changes. Those descriptive interval counts are not
interchangeable with the Holm-corrected permutation tests.

Neither newspaper has an unexpectedly high reversal \emph{count}
under this conditional reference. The 29 adjusted \emph{Imvo}
estimator gaps answer a different question: nine also had original
sign-reversal warnings. Table~\ref{tab:sup-examples} shows several
of those descriptive disagreements, using the same patterns and
original corpus. For \texttt{ants}, the maximum gap has
Holm-adjusted $p=0.024$, whereas the \emph{mass-weighted temporal}
test has adjusted $p=1$. The gap is evidence of estimator sensitivity,
not a pair of proven historical trends.

\begin{table}[H]
\centering\footnotesize
\renewcommand{\arraystretch}{1.15}
\begin{tabular*}{\columnwidth}{@{\extracolsep{\fill}}lrrrr@{}}
\hline
Fragment & Mass & Page & Issue & Month\\
\hline
\texttt{ants} & $-0.23$ & $+4.80$ & $-0.03$ & $+0.15$\\
\texttt{ntsu} & $-2.03$ & $+0.81$ & $-1.79$ & $-1.65$\\
\texttt{tsha} & $+0.40$ & $-2.23$ & $+0.20$ & $+0.12$\\
\texttt{teto} & $-0.36$ & $+2.21$ & $-0.12$ & $-0.29$\\
\hline
\end{tabular*}
\caption{Illustrative \emph{Imvo} 1888-minus-1885 rate
changes per 10,000 eligible four-gram positions. Patterns are
surface fragments, not validated isiXhosa words or meanings;
individual sign changes are not adjusted discoveries.}
\label{tab:sup-examples}
\end{table}

\clearpage
\subsection*{S4. Archive structure, repeated text, and validity}
Figure~\ref{fig:sup-page} shows how eligible page text varies with
position. \emph{Imvo} page 2 has 4.24 times the within-word four-gram
opportunities of page 1 in 1885, and 3.69 times in 1888; the
\emph{Isigidimi} profile is considerably flatter. The same pattern
can therefore carry different weight in character-mass and equal-page
comparisons. Page position is not an annotated genre or evidence of
editorial motives.

\begin{figure}[H]
\centering
\includegraphics[width=.99\columnwidth]{supp_page_position_compact.pdf}
\caption{Mean eligible four-gram positions per page relative to
position 1 in each newspaper and year, using the original XH-assigned
archive text. Positions 1--4 are present in all four source/year
combinations. This is exposure, not a rate of a selected pattern.}
\label{fig:sup-page}
\end{figure}

The follow-up masks words covered by exact 13-word sequences
repeated across \emph{different} issues within the same newspaper.
Table~\ref{tab:sup-cover} shows the text affected. After masking,
ten original \emph{Imvo} reversal warnings and one
\emph{Isigidimi} warning disappear;
this changes the eligible text and is a sensitivity analysis, not an
estimate of the causal effect of ``removing duplication.'' Among the
32 originally flagged \emph{Imvo} patterns, seven can change their
mass-weighted year-difference sign when one complete issue is omitted;
four of six flagged \emph{Isigidimi} patterns do so.

\begin{table}[H]
\centering\footnotesize
\renewcommand{\arraystretch}{1.15}
\begin{tabular*}{\columnwidth}{@{\extracolsep{\fill}}lrrr@{}}
\hline
Newspaper / year & Pages & Issues & Masked words\\
\hline
\emph{Imvo} 1885 & 195 & 49 & 20.0\%\\
\emph{Imvo} 1888 & 208 & 52 & 21.9\%\\
\emph{Isigidimi} 1885 & 95 & 12 & 2.4\%\\
\emph{Isigidimi} 1888 & 96 & 12 & 1.3\%\\
\hline
\end{tabular*}
\caption{Original eligible page support and fraction of words
covered by repeated passages in each newspaper and year.
The mask also removes some pages from the sensitivity analysis.}
\label{tab:sup-cover}
\end{table}

The tokenizer experiment also has disjoint date-based splits but not
perfectly novel evaluation \emph{content}: 514 of 2,493 held-out
page-view excerpts (20.6\%) contain an exact 13-word sequence also
seen in training-era text. Restricting to 99 nonoverlapping 1888 English-view excerpts yields a
post-hoc BPE bigram allocation effect of $-0.02874$ BPC, compared with
$-0.03722$ on full support. The direction remains, but the restriction
changes the evaluation population, and shorter or paraphrased overlaps
remain possible.

ArchiveLens uses the original text views, not fitted tokenizer
outputs. The two studies share a question about measurement, not a
measured causal pathway. Both use surviving transcriptions from one
archive, with no experimental control of editorial composition,
preservation or OCR/transcription effects. The historical years
had been examined before the final follow-ups, and permutation
exchangeability is a conditional reference rather than a credible
assumption of identical historical eras. The findings support
measurement- and sampling-sensitivity diagnostics, not claims about
historical authors' beliefs, semantic change, social discrimination,
or downstream neural-model fairness.

\noindent\textbf{Generative-AI assistance.}
ChatGPT assisted experiment planning, analysis/code checks, figure preparation,
and report drafting/editing; NotebookLM and Gemini assisted literature
exploration and citations, as declared. The authors remain responsible
for validating the code, numerical evidence and references.

\noindent\textbf{Code and reproduction.}
The recorded model and ArchiveLens results are verified by the
repository's Python 3.13 checks; the full tokenizer retraining requires
the original source ZIP and is not run in routine CI. The source archive
SHA-256 is \texttt{375c5864...be321d5} (full value in
\texttt{README.txt}). Hardware and runtime measurements were not
retained consistently, so none are estimated here. Archive redistribution
rights require separate confirmation.
The tables can be checked in
\path{results/segmentation.csv}, \path{results/bigram.csv}, and
\path{results/repetition.csv}. The 120-pattern estimands and their
adjusted calibration tests are in \path{results/archivelens/}; numerical inputs for
Fig.~\ref{fig:sup-page} are generated directly from the original ZIP
by \path{paper/analysis/make_supp_figure.py}; the independently
recounted inputs are in \path{paper/analysis/supp_page_position_summary.csv}.
Reproduction commands appear in the required \path{README.txt}.
 AFTER the
% references and completed, author-approved NeurIPS 2024 questionnaire,
% as specified by the course brief. Requires graphicx (in the CCN manuscript).
% This is not a replacement for the required Faculty AI ethics statement.
% This file is separate until the NeurIPS answers are author-approved.
% It must not be appended before that required questionnaire.
\clearpage
\section*{Supplementary Material}

\subsection*{S1. Sampling design and predictive evaluation}
\noindent\textbf{Corpus and separation.}
The tokenizer experiment uses the \emph{Imvo Zabantsundu} pages in the supplied
language-view archive. It trains on 1885 and January--September 1887
(342 pages, 86 publication issues), selects the count-bigram smoothing
parameter on October--December 1887 (48 pages, 12 issues), and evaluates
on 1888 (52 issues; 134 English-assigned and 208 isiXhosa-assigned
eligible page-view excerpts). The 1889--1894 material is a dependent sensitivity check (306 issues):
those years had been inspected in earlier exploratory work. Training and evaluation
issue identifiers are disjoint, but identical passages can recur in
different newspaper issues. Eligible evaluation page views contribute their first 128 alphabetic
words when at least 40 are present. English/isiXhosa views are archive
assignments, not manually adjudicated word-level languages.

\noindent\textbf{Four sampling arms.}
For eligible view $v$ of page $i$, let $n_{iv}$ be the number of ASCII
alphabetic words, $n_i=\sum_v n_{iv}$, and $N$ the number of eligible
pages. The probability of selecting a language view is either the
page-average share or the pooled word share:
\begin{equation}
\begin{aligned}
p_P(v)&=\frac{1}{N}\sum_i\frac{n_{iv}}{n_i},\\
p_W(v)&=\frac{\sum_i n_{iv}}{\sum_i n_i}.
\end{aligned}
\end{equation}
Conditional on the view, the page is drawn with probability
$q_P(i\mid v)\propto n_{iv}/n_i$ or
$q_W(i\mid v)\propto n_{iv}$, normalized over its eligible pages.
In the training pool, the expected English-view draw share is
23.55\% for $P$ and 32.56\% for $W$; realized character shares can differ
because sampled windows vary in length.
The cross of allocation ($P/W$) and within-view page selection ($P/W$)
produces $PP$, $PW$, $WP$ and $WW$, respectively; changing only the first
letter isolates the allocation factor. Each draw selects a uniformly
random eligible 32-word within-page window. The study uses lowercase
\texttt{[a-z]+} words, no stemming or modernization, and a training
budget of $500{,}000\pm30$ alphabetic characters per tokenizer.

Both BPE and Unigram use a fixed 2,000-piece vocabulary, identity
SentencePiece normalization, character coverage 1.0, single-threaded
fitting, and five sampling seeds (441--445) per arm:
$2\times4\times5=40$ main tokenizers.
The vocabulary size and tokenizer families were chosen in earlier
exploration and held fixed for this study, not selected on the 1888 set. The predictive model is a
\emph{count bigram}, not a neural language model. Each tokenizer is
applied to the same underlying 545 training page-view texts and held-out
raw excerpts; only the induced token sequence changes.

\noindent\textbf{Measures and model selection.}
Segmentation cost is $E=100P/C$, where $P$ is the number of SentencePiece
tokens and $C$ is the number of alphabetic source characters. For a token
$t$ and predecessor $h$, with token vocabulary $V$ and total training
token count $T$, the smoothed bigram probability is
\begin{equation}
\begin{aligned}
u(t)&=\frac{c(t)+0.1}{T+0.1|V|},\\
\Pr(t\mid h)&=\frac{c(h,t)+\alpha u(t)}{c(h)+\alpha}.
\end{aligned}
\end{equation}
The start-of-document predecessor is handled as a separate context.
Held-out bits per alphabetic character (BPC) is the negative base-two
log likelihood divided by $C$. The globally shared $\alpha$ was
chosen from $\{1,10,100\}$ on 344 validation windows screened for
exact training 13-word overlaps (13,760 model-window scores), averaging
view-level BPC across tokenizer families, arms and seeds. The respective validation
means were 3.1246, 2.8279 and 2.7329 BPC, selecting $\alpha=100$ before
the 1888 comparison.

For either response $Y\in\{E,\mathrm{BPC}\}$, the factorial
allocation and within-view contrasts are
\begin{equation}
\begin{aligned}
\Delta_A&=\tfrac12[(Y_{WP}-Y_{PP})+(Y_{WW}-Y_{PW})],\\
\Delta_B&=\tfrac12[(Y_{PW}-Y_{PP})+(Y_{WW}-Y_{WP})].
\end{aligned}
\end{equation}
The reported effects are means over five fitted seeds; 95\% percentile
intervals use 5,000 paired resamples of \emph{entire publication issues}
(52 issues for 1888), retaining all pages and arms of a sampled issue
together. These intervals describe uncertainty over the observed issue clusters,
conditional on the fixed trained seeds and archive; five seeds are not
five independent historical samples.

\subsection*{S2. Full tokenizer results and falsification checks}
The main paper emphasizes BPE English and the language-view trade-off.
Tables~\ref{tab:sup-seg}--\ref{tab:sup-bpc} give the complete primary
1888 allocation contrast for both models and both archive views.
Negative values indicate lower cost under word-like view allocation.

\begin{table}[ht]
\centering\footnotesize
\renewcommand{\arraystretch}{1.13}
\begin{tabular*}{\columnwidth}{@{\extracolsep{\fill}}lrr@{}}
\hline
Model/view & Effect & 95\% interval\\
\hline
BPE / EN & $-1.878$ & $[-1.972,-1.782]$\\
BPE / XH & $+0.453$ & $[+0.407,+0.497]$\\
Unigram / EN & $-1.795$ & $[-1.908,-1.671]$\\
Unigram / XH & $+0.555$ & $[+0.513,+0.600]$\\
\hline
\end{tabular*}
\caption{1888 allocation effect $\Delta_A$: subword pieces per
100 alphabetic characters, with issue-bootstrap 95\% intervals.}
\label{tab:sup-seg}
\end{table}

\begin{table}[ht]
\centering\footnotesize
\renewcommand{\arraystretch}{1.13}
\begin{tabular*}{\columnwidth}{@{\extracolsep{\fill}}lrr@{}}
\hline
Model/view & Effect & 95\% interval\\
\hline
BPE / EN & $-0.03722$ & $[-0.04236,-0.03212]$\\
BPE / XH & $-0.00275$ & $[-0.00523,-0.00014]$\\
Unigram / EN & $-0.03834$ & $[-0.04366,-0.03318]$\\
Unigram / XH & $+0.01048$ & $[+0.00805,+0.01286]$\\
\hline
\end{tabular*}
\caption{1888 $\Delta_A$: count-bigram bits per alphabetic character.
Held-out source texts are matched across arms.}
\label{tab:sup-bpc}
\end{table}

BPE's XH bigram change is small and has mixed signs across seeds
(three negative, two positive), unlike its five-of-five positive
segmentation effect. Unigram has higher XH segmentation \emph{and}
higher XH bigram loss. Thus a tokenizer can encode the archive more
compactly in one view without consistently improving predictive loss
in the other. In English-view BPE, the separately estimated
within-view selection effect on BPC is $+0.01938$
(95\% interval $[+0.01146,+0.02769]$), opposing the allocation benefit.

For the repetition-control experiment, 20 additional BPE tokenizers
were trained under $PP/WW$, five seeds and two matched-support
conditions. The 1888 English $E_{PP}-E_{WW}$ contrast after filtering
training-only repeated 13-word sequences was $+2.389$
$[+2.182,+2.590]$ pieces per 100 characters; an unfiltered model on
the same page-view support gave $+2.097$ $[+1.905,+2.280]$.
Their paired difference is $+0.292$ $[+0.163,+0.421]$.
The trade-off survives this exact reprint filter;
the filter cannot rule out shorter, paraphrased or otherwise shared
historical passages.

A counterexample concerns the \emph{whole-policy}, rather than
allocation-only, contrast. On dependent 1889--1894 English data,
BPE $\mathrm{BPC}_{WW}-\mathrm{BPC}_{PP}=+0.00450$; the issue-bootstrap
interval is $[+0.00176,+0.00728]$, but the six-year resampling
interval $[-0.00268,+0.01305]$ includes zero. This does not support
a reliable improvement from changing the entire sampling policy.

\subsection*{S3. ArchiveLens estimands and statistical calibration}
ArchiveLens uses the archive's XH-assigned views from both
\emph{Imvo Zabantsundu} and \emph{Isigidimi sama-Xosa}.
It fixes the 120 most frequent within-word four-character sequences
by the mean of the two newspapers' normalized \emph{1887} rates;
there is no selection from 1885/1888 outcome differences.
A page is eligible if it has at least 40 alphabetic words.
The 1885/1888 support comprises 195/208 pages (49/52 issues)
for \emph{Imvo} and 95/96 pages (12/12 issues) for \emph{Isigidimi}.

For pattern $g$ on page $i$, let $c_{ig}$ be its occurrence count and
$d_i=\sum_{w\in i}\max(0,|w|-3)$ the number of eligible within-word
four-gram positions. The \emph{mass} rate pools all counts and
positions, while \emph{page} gives each eligible page the same weight.
The \emph{issue} rate first pools pages within an issue, then averages
issue rates. The \emph{month} rate averages issue rates within each
calendar month before weighting the 12 shared calendar months equally.
All rates are multiplied by $10^4$. Each reported historical contrast
is the 1888 rate minus the 1885 rate. These target different exposure units rather than a single,
unambiguously correct historical population.

A 3,000-draw nested month/issue bootstrap produces descriptive
95\% intervals, not multiplicity-adjusted discovery claims.
Separately, 4,999 year-label permutations within calendar months move
whole publication issues, preserve month/year issue counts and retain
page groups. For each pattern, one test examines
$|\Delta_{g,\mathrm{mass}}|$ (temporal change), while another examines
$\max_{w\ne\mathrm{mass}}|\Delta_{g,w}-\Delta_{g,\mathrm{mass}}|$
(estimator sensitivity). Each family of 120 pattern tests within a
newspaper receives its own Holm correction. A third benchmark examines
the \emph{number} of sign reversals, not the reliability of each
individual reversal warning.

\begin{table}[H]
\centering\footnotesize
\renewcommand{\arraystretch}{1.13}
\begin{tabular*}{\columnwidth}{@{\extracolsep{\fill}}lrr@{}}
\hline
ArchiveLens result & Imvo & Isigidimi\\
\hline
Reversal warnings & 23/120 & 4/120\\
Null mean & 29.02 & 8.30\\
Central 95\% null interval & 19--40 & 3--16\\
Upper-tail $p$ & 0.8882 & 0.9486\\
Holm temporal effects & 26 & 0\\
Holm estimator gaps & 29 & 0\\
\hline
\end{tabular*}
\caption{Reversal-count calibration and separately Holm-adjusted year and estimator-gap tests.}
\label{tab:sup-arch}
\end{table}

For context, the earlier \emph{unadjusted} bootstrap intervals
excluded zero for 38 \emph{Imvo} and 21 \emph{Isigidimi}
mass-weighted year changes. Those descriptive interval counts are not
interchangeable with the Holm-corrected permutation tests.

Neither newspaper has an unexpectedly high reversal \emph{count}
under this conditional reference. The 29 adjusted \emph{Imvo}
estimator gaps answer a different question: nine also had original
sign-reversal warnings. Table~\ref{tab:sup-examples} shows several
of those descriptive disagreements, using the same patterns and
original corpus. For \texttt{ants}, the maximum gap has
Holm-adjusted $p=0.024$, whereas the \emph{mass-weighted temporal}
test has adjusted $p=1$. The gap is evidence of estimator sensitivity,
not a pair of proven historical trends.

\begin{table}[H]
\centering\footnotesize
\renewcommand{\arraystretch}{1.15}
\begin{tabular*}{\columnwidth}{@{\extracolsep{\fill}}lrrrr@{}}
\hline
Fragment & Mass & Page & Issue & Month\\
\hline
\texttt{ants} & $-0.23$ & $+4.80$ & $-0.03$ & $+0.15$\\
\texttt{ntsu} & $-2.03$ & $+0.81$ & $-1.79$ & $-1.65$\\
\texttt{tsha} & $+0.40$ & $-2.23$ & $+0.20$ & $+0.12$\\
\texttt{teto} & $-0.36$ & $+2.21$ & $-0.12$ & $-0.29$\\
\hline
\end{tabular*}
\caption{Illustrative \emph{Imvo} 1888-minus-1885 rate
changes per 10,000 eligible four-gram positions. Patterns are
surface fragments, not validated isiXhosa words or meanings;
individual sign changes are not adjusted discoveries.}
\label{tab:sup-examples}
\end{table}

\clearpage
\subsection*{S4. Archive structure, repeated text, and validity}
Figure~\ref{fig:sup-page} shows how eligible page text varies with
position. \emph{Imvo} page 2 has 4.24 times the within-word four-gram
opportunities of page 1 in 1885, and 3.69 times in 1888; the
\emph{Isigidimi} profile is considerably flatter. The same pattern
can therefore carry different weight in character-mass and equal-page
comparisons. Page position is not an annotated genre or evidence of
editorial motives.

\begin{figure}[H]
\centering
\includegraphics[width=.99\columnwidth]{supp_page_position_compact.pdf}
\caption{Mean eligible four-gram positions per page relative to
position 1 in each newspaper and year, using the original XH-assigned
archive text. Positions 1--4 are present in all four source/year
combinations. This is exposure, not a rate of a selected pattern.}
\label{fig:sup-page}
\end{figure}

The follow-up masks words covered by exact 13-word sequences
repeated across \emph{different} issues within the same newspaper.
Table~\ref{tab:sup-cover} shows the text affected. After masking,
ten original \emph{Imvo} reversal warnings and one
\emph{Isigidimi} warning disappear;
this changes the eligible text and is a sensitivity analysis, not an
estimate of the causal effect of ``removing duplication.'' Among the
32 originally flagged \emph{Imvo} patterns, seven can change their
mass-weighted year-difference sign when one complete issue is omitted;
four of six flagged \emph{Isigidimi} patterns do so.

\begin{table}[H]
\centering\footnotesize
\renewcommand{\arraystretch}{1.15}
\begin{tabular*}{\columnwidth}{@{\extracolsep{\fill}}lrrr@{}}
\hline
Newspaper / year & Pages & Issues & Masked words\\
\hline
\emph{Imvo} 1885 & 195 & 49 & 20.0\%\\
\emph{Imvo} 1888 & 208 & 52 & 21.9\%\\
\emph{Isigidimi} 1885 & 95 & 12 & 2.4\%\\
\emph{Isigidimi} 1888 & 96 & 12 & 1.3\%\\
\hline
\end{tabular*}
\caption{Original eligible page support and fraction of words
covered by repeated passages in each newspaper and year.
The mask also removes some pages from the sensitivity analysis.}
\label{tab:sup-cover}
\end{table}

The tokenizer experiment also has disjoint date-based splits but not
perfectly novel evaluation \emph{content}: 514 of 2,493 held-out
page-view excerpts (20.6\%) contain an exact 13-word sequence also
seen in training-era text. Restricting to 99 nonoverlapping 1888 English-view excerpts yields a
post-hoc BPE bigram allocation effect of $-0.02874$ BPC, compared with
$-0.03722$ on full support. The direction remains, but the restriction
changes the evaluation population, and shorter or paraphrased overlaps
remain possible.

ArchiveLens uses the original text views, not fitted tokenizer
outputs. The two studies share a question about measurement, not a
measured causal pathway. Both use surviving transcriptions from one
archive, with no experimental control of editorial composition,
preservation or OCR/transcription effects. The historical years
had been examined before the final follow-ups, and permutation
exchangeability is a conditional reference rather than a credible
assumption of identical historical eras. The findings support
measurement- and sampling-sensitivity diagnostics, not claims about
historical authors' beliefs, semantic change, social discrimination,
or downstream neural-model fairness.

\noindent\textbf{Generative-AI assistance.}
ChatGPT assisted experiment planning, analysis/code checks, figure preparation,
and report drafting/editing; NotebookLM and Gemini assisted literature
exploration and citations, as declared. The authors remain responsible
for validating the code, numerical evidence and references.

\noindent\textbf{Code and reproduction.}
The recorded model and ArchiveLens results are verified by the
repository's Python 3.13 checks; the full tokenizer retraining requires
the original source ZIP and is not run in routine CI. The source archive
SHA-256 is \texttt{375c5864...be321d5} (full value in
\texttt{README.txt}). Hardware and runtime measurements were not
retained consistently, so none are estimated here. Archive redistribution
rights require separate confirmation.
The tables can be checked in
\path{results/segmentation.csv}, \path{results/bigram.csv}, and
\path{results/repetition.csv}. The 120-pattern estimands and their
adjusted calibration tests are in \path{results/archivelens/}; numerical inputs for
Fig.~\ref{fig:sup-page} are generated directly from the original ZIP
by \path{paper/analysis/make_supp_figure.py}; the independently
recounted inputs are in \path{paper/analysis/supp_page_position_summary.csv}.
Reproduction commands appear in the required \path{README.txt}.
